\chapter{Terminal Configuration Schema}\label{app:terminal-schema}

This appendix documents the JSON schema of the Terminal Configuration, which defines the behaviour of the simulated terminal in the \pcbctf{} platform.

% ============================================================================
\section{Root Structure}\label{sec:app-tc-root}
% ============================================================================

\noindent\begin{minipage}{\linewidth}
\begin{lstlisting}[language={}, caption={Top-level TerminalConfig schema.}]
{
  "metadata": {
    "version": "1.0.0",
    "name": "WR841N Terminal",
    "description": "UART console simulation"
  },
  "tabs": [ ... ],
  "flags": {
    "parts": [ ... ],
    "completeFlag": "flag{b00t_r00t_h4sh_l34k}",
    "showProgress": true
  },
  "globalCommands": { ... }
}
\end{lstlisting}
\end{minipage}

% ============================================================================
\section{Tab Configuration}\label{sec:app-tc-tab}
% ============================================================================

Each tab defines an independent terminal environment with its own filesystem, commands, and optional boot sequence.

\noindent\begin{minipage}{\linewidth}
\begin{lstlisting}[language={}, caption={TabConfig structure.}]
{
  "id": "uart",
  "name": "UART Console",
  "filesystem": { "tree": { ... } },
  "commands": { ... },
  "bootSequence": {
    "stages": [ ... ],
    "initialStage": "connecting"
  },
  "initialPath": "/",
  "defaultConstraints": { ... }
}
\end{lstlisting}
\end{minipage}

% ============================================================================
\section{Command Definition}\label{sec:app-tc-cmd}
% ============================================================================

Commands are the core building blocks of the terminal simulation. Each command specifies a handler type, optional constraints, an output definition, and optional side effects.

\noindent\begin{minipage}{\linewidth}
\begin{lstlisting}[language={}, caption={CommandDefinition with lookup output.}]
{
  "description": "Extract printable strings",
  "handler": "custom",
  "constraints": {
    "state": { "bootStage": "shell" }
  },
  "output": {
    "type": "lookup",
    "argIndex": 0,
    "matchType": "contains",
    "table": {
      "/dev/mtdblock3": [
        "admin_pass=sohoadmin",
        "wifi_password=myWiFi123"
      ],
      "/usr/bin/httpd": [
        "execFormatCmd(%s)"
      ]
    },
    "default": { "type": "static", "lines": ["..."] }
  },
  "sideEffects": {
    "unlockFlags": [
      {
        "id": "leak",
        "condition": {
          "type": "argument",
          "contains": "mtdblock3"
        }
      }
    ]
  }
}
\end{lstlisting}
\end{minipage}

\subsection{Handler Types}\label{subsec:app-handlers}

\begin{table}[H]
\centering
\caption{Command handler types.}
\label{tab:handlers}
\begin{tabularx}{\textwidth}{lX}
\toprule
\textbf{Type} & \textbf{Description} \\
\midrule
\texttt{builtin}  & Delegates to a built-in handler (ls, cd, cat, pwd, grep, find, clear) \\
\texttt{custom}   & Routes to the configured output definition \\
\texttt{dynamic}  & Calls a registered custom function by name \\
\texttt{script}   & Calls a registered custom function with a different signature \\
\bottomrule
\end{tabularx}
\end{table}

\subsection{Output Types}\label{subsec:app-outputs}

\begin{table}[H]
\centering
\caption{Terminal output types.}
\label{tab:outputs}
\begin{tabularx}{\textwidth}{lX}
\toprule
\textbf{Type} & \textbf{Description} \\
\midrule
\texttt{static}      & Fixed array of output lines \\
\texttt{conditional}  & If-then-else chain evaluated against the execution context \\
\texttt{template}     & String with \texttt{\{\{variable\}\}} placeholders \\
\texttt{lookup}       & Table indexed by a command argument \\
\texttt{dynamic}      & Custom function invocation \\
\texttt{script}       & Custom function with alternate signature \\
\bottomrule
\end{tabularx}
\end{table}

\subsection{Constraint Categories}\label{subsec:app-constraints}

\begin{table}[H]
\centering
\caption{Command constraint categories and evaluation order.}
\label{tab:constraints}
\begin{tabularx}{\textwidth}{clX}
\toprule
\textbf{Order} & \textbf{Category} & \textbf{Description} \\
\midrule
1 & \texttt{path}          & Restricts the directories from which the command can be invoked \\
2 & \texttt{permissions}   & Requires specific user identity (e.g., root) \\
3 & \texttt{prerequisites} & Requires prior commands, files, or flags \\
4 & \texttt{arguments}     & Validates argument count and patterns \\
5 & \texttt{state}         & Requires specific boot stage or state variables \\
\bottomrule
\end{tabularx}
\end{table}

% ============================================================================
\section{Filesystem Structure}\label{sec:app-tc-fs}
% ============================================================================

The filesystem can be specified in a nested tree format (authoring-friendly) that is automatically flattened by the configuration loader:

\noindent\begin{minipage}{\linewidth}
\begin{lstlisting}[language={}, caption={Filesystem tree notation.}]
{
  "tree": {
    "etc": {
      "passwd": "root:x:0:0:root:/root:/bin/sh",
      "shadow": "root:$1$abc:18000:0:99999:7:::"
    },
    "proc": {
      "cpuinfo": "system type: QCA953x ..."
    },
    "dev": {}
  }
}
\end{lstlisting}
\end{minipage}

String values represent file contents; empty objects represent empty directories. The loader produces two parallel dictionaries: \texttt{directories} (path $\rightarrow$ entry names) and \texttt{files} (path $\rightarrow$ content or metadata).

% ============================================================================
\section{Boot Sequence}\label{sec:app-tc-boot}
% ============================================================================

Boot sequences define a chain of stages that simulate the device startup process:

\noindent\begin{minipage}{\linewidth}
\begin{lstlisting}[language={}, caption={Boot stage example.}]
{
  "id": "uboot_wait",
  "name": "U-Boot Autoboot",
  "lines": ["Hit any key to stop autoboot: "],
  "duration": 3000,
  "nextStage": "kernel_boot",
  "prompt": "",
  "autoProgressTimeout": 5000
}
\end{lstlisting}
\end{minipage}

Three stage identifiers receive special handling:
\begin{itemize}
    \item \texttt{uboot\_wait}: Displays a countdown and advances to the next stage if the user presses a key, or auto-progresses after the timeout.
    \item \texttt{login}: Prompts for a username (displayed as typed).
    \item \texttt{password}: Prompts for a password (displayed as asterisks).
\end{itemize}

% ============================================================================
\section{Flag System}\label{sec:app-tc-flags}
% ============================================================================

\noindent\begin{minipage}{\linewidth}
\begin{lstlisting}[language={}, caption={Flag system definition.}]
{
  "parts": [
    { "id": "boot", "part": "b00t",
      "description": "Boot analysis", "hint": "..." },
    { "id": "root", "part": "_r00t",
      "description": "Root access", "hint": "..." }
  ],
  "completeFlag": "flag{b00t_r00t_h4sh_l34k_1nj3ct_sh3ll}",
  "showProgress": true
}
\end{lstlisting}
\end{minipage}

Flag parts are unlocked through command-level side effects. The terminal UI displays discovered fragments and masks pending ones with \texttt{?} characters.
